%! Piecewise \& Step Functions — Unit 2, Lesson 2.4 — Warm-Up (Answer Key)
\documentclass[10pt]{article}
\usepackage{algebra2-article}
\usepackage{algebra2-key}   % pulls in -boxes; do NOT also load -boxes
\newcommand{\TallMath}[1]{$\displaystyle #1\rule[-1.4em]{0pt}{3.2em}$}

\begin{document}

\pageheader{Unit 2, Lesson 2.4}{Warm-Up --- Answer Key}
\namedateperiod

\vspace{0.10in}

\begin{notesbox}{Getting Ready: Rules, Two Pieces, and Endpoints}
\small Today a function can follow more than one rule. These three quick reviews are the tools you'll
need.

\vspace{4pt}
\textbf{1. Evaluate the rule.} For $g(x)=2x-1$, find each output.
\begin{center}
\renewcommand{\arraystretch}{1.25}
\begin{tabular}{|c|c|c|c|}
\hline
$x$ & $-2$ & $0$ & $3$ \\ \hline
$g(x)=2x-1$ & ${\color{keyred}\mathbf{-5}}$ & ${\color{keyred}\mathbf{-1}}$ & ${\color{keyred}\mathbf{5}}$ \\ \hline
\end{tabular}
\end{center}
And for $h(x)=|x|$: \quad $h(-4)={\color{keyred}\mathbf{4}}$, \quad $h(0)={\color{keyred}\mathbf{0}}$.

\vspace{6pt}
\textbf{2. The V is two lines (from Lesson 2.3).} The parent $y=|x|$ is made of two straight lines
pasted at the vertex. Fill in each rule.
\[ y=|x| \text{ behaves like } y={\color{keyred}\mathbf{x}} \text{ when } x\ge 0, \text{ and like }
   y={\color{keyred}\mathbf{-x}} \text{ when } x<0. \]

\vspace{4pt}
\textbf{3. Endpoints and inequalities.} A number line uses a \emph{closed} dot $\bullet$ to
\emph{include} a point and an \emph{open} dot $\circ$ to \emph{exclude} it.
\begin{itemize}[leftmargin=1.6em, itemsep=3pt]
  \item Is $x=1$ a solution of $x<1$? \ans{no} \quad Is $x=1$ a solution of $x\le 1$? \ans{yes}
  \item To show ``$x\ge 2$'' on a number line, the dot at $2$ should be \ans{closed} (open / closed).
\end{itemize}
\end{notesbox}

\begin{teachernote}
Item~1 is the mechanical core: each ``piece'' of a piecewise function is just a rule you evaluate, so
students must be fluent plugging in --- including into $|x|$. Item~2 is the whole bridge from
Lesson~2.3: the V is literally two linear rules, $y=x$ and $y=-x$, joined at the vertex; that is the
first piecewise function students already know. Item~3 is the make-or-break idea for today: $x=1$
satisfies $x\le 1$ but not $x<1$, which is exactly how a boundary point is assigned to one piece (and
drawn with a closed vs.\ open dot). Debrief item~3 aloud before the notes.
\end{teachernote}

\end{document}
